\section{Pigeonhole critical assignments and restriction helpers}

\begin{annexitem}{annex:p7.haken.critical-count-routing}{thm:p07-haken-critical-count-routing}{Count, false inputs and every pivot routing case}
\begin{proof}
There are n+1 choices for u and n! bijections of the other
n labeled pigeons to the n labeled holes. Distinct choices
give different full assignments: the all-false row identifies
u, and each other row identifies its image. Conversely, every
defined critical assignment is one of these choices. Thus the
count is \((n+1)n!\). An assigned row has one true literal,
and row u has none. For each collision clause, two distinct
pigeons cannot both have its hole as their assigned image,
so at least one of its negative literals is true. All inputs
except \(A_u\) are therefore satisfied.

Consider a resolution node with parents E and F, pivot literal
L in E and its complement in F, and resolvent
\((E\setminus\{L\})\cup(F\setminus\{\overline L\})\).
If an assignment falsifies the resolvent, all its nonpivot
literals are false. If L is false, E is false and is the
selected parent; if L is true, its complement is false and
F is false and is selected. The other parent has its true
pivot literal and is satisfied. Thus the choice is unique
for the fixed annotated inference. A tautological node has
no falsifying assignment and never lies on such a route.

The empty clause is false under every assignment, so routing
starts there. Parent references point to earlier nodes; a
finite strictly decreasing sequence of node indices terminates
at an input. That input is false by induction along the route,
and for a critical assignment only \(A_u\) can be false.
Shared parents do not affect either its truth or the unique
pivot choice. This proves routing for DAGs without replacing
the distinct assignment count by a path count.
\end{proof}
\end{annexitem}

\begin{annexitem}{annex:p7.haken.restricted-family}{thm:p07-haken-restricted-family}{All residual input cases and the exact conditional count}
\begin{proof}
An assigned row has its designated true literal and is discarded.
An unassigned row has false literals precisely in J and remains
the disjunction of its n-r unoccupied columns. For collisions
between two unassigned pigeons, an occupied column has both
variables false, so that collision clause is satisfied and
discarded; an unoccupied column is unset and its collision
clause remains unchanged. For a collision with an assigned
pigeon i, if its column is f(i), the other pigeon's variable
in that occupied column is false, satisfying the negative
literal. In every other column the assigned pigeon's variable
itself is false, again satisfying the clause. With two assigned
pigeons, their injectively distinct images ensure at least
one false variable in every common column. All such collisions
are discarded. These exhaust the input cases. The remaining
rows and unoccupied columns consequently form exactly the
claimed smaller nonfunctional CNF, with no additional clauses.

For r=0, no variable is set and the family is unchanged.
For r=n, only one pigeon remains and its residual row has
no literals, giving the declared empty-clause instance.
Every assigned row's designated true literal prevents that
pigeon from being the unmatched one in a compatible critical
assignment. Thus u has n+1-r possible choices outside I.
After its choice, the remaining n-r pigeons must map bijectively
to the n-r unoccupied holes; there are \((n-r)!\) choices.
Each choice extends all of \(\rho_f\), and every compatible
critical assignment has exactly this form. Removing the fixed
rows and occupied columns is a bijection with the residual
critical family. Its count is therefore
\((n+1-r)(n-r)!\), including the endpoint value one when
r=n. Dividing by the full critical count proves the exact
fraction, without an independence assumption.
\end{proof}
\end{annexitem}

\begin{annexitem}{annex:p7.haken.resolution-restriction}{thm:p07-haken-resolution-restriction}{Subclause tracking rather than an added weakening inference}
\begin{proof}
Process the original proof in its topological order. For every
node C whose restriction is not the satisfied marker, maintain
a derived clause \(C'\subseteq C\upharpoonright\rho\).
A surviving input has its restricted input itself as \(C'\).
For a resolution node C with pivot literal L and parents E,F,
assume \(C\upharpoonright\rho\ne\topv\).
All nonpivot parent literals occur in C, so none is set true.

If L is set true, F is not satisfied: its complementary pivot
is false and every remaining literal is unsatisfied. Moreover
\(F\upharpoonright\rho\subseteq C\upharpoonright\rho\),
because its false pivot is deleted. Reuse the tracked F' for
C'. If L is set false, the symmetric argument reuses E'.
This handles both assigned-pivot cases without producing a
new clause or using a deleted satisfied parent.

If the pivot is unset, both parent restrictions are unsatisfied
and have tracked E',F'. If E' lacks L, then every literal in
E' is among C's nonpivot residual literals, so reuse E'.
If F' lacks the complement, reuse F' by the same reasoning.
Otherwise resolve E' and F' on that pivot. The sheet's literal-set
union gives a subclause of \(C\upharpoonright\rho\), so this
one actual resolution step supplies C'. This includes clauses
where pruning has made the pivot unnecessary; no inference from
a smaller clause to a larger one was assumed.

At the original empty-clause node its restriction is the empty
clause, whose only subclause is itself. The tracked final node
is therefore an actual empty clause derived from surviving
restricted inputs. At most one input or resolution node was
created per original surviving node; reuse creates no additional
occurrence. The new refutation has no more nodes and can retain
sharing. For an original tree, perform the same construction
on its occurrences separately: pruning selects a parent subtree
or combines the two restricted subtrees. No crossing edge or
new duplication is created, so its output is again a tree with
no more occurrences. This proves both style-specific claims.
\end{proof}
\end{annexitem}

\begin{annexitem}{annex:p7.haken.charging}{thm:p07-haken-charging}{Finite counting with explicit barrier and capacity hypotheses}
\begin{proof}
Every critical route has a first node in S by the barrier
hypothesis and finiteness of its route. Charge it to that
unique node. The resulting fibers partition \(\mathcal C_n\).
If their individual counts are \(a_v\), then
\[
 (n+1)n!=\sum_{v\in S}a_v\leq b\,\#S.
\]
Divide by positive b and use that \(\#S\) is an integer.
If b is at most \(2^{-\beta n}(n+1)n!\), cancel the
positive factorial count to obtain \(\#S\geq2^{\beta n}\).
These algebraic steps use both hypotheses; they construct
neither a special barrier nor its capacity estimate.

For the row-input barrier, exactly the n! assignments with a
given unmatched u terminate at row u, so its capacity is n!
and the inequality yields n+1. For the singleton empty-clause
barrier, every assignment starts at its only node, so the
capacity is \((n+1)n!\) and the inequality yields one.
Both examples satisfy the formal charging hypotheses but
produce no exponential estimate. They exhibit the exact
additional content an exponential bottleneck lemma must supply.
\end{proof}
\end{annexitem}

\begin{annexitem}{annex:p7.haken.small-routes}{thm:p07-haken-small-routes}{Every one of the six assignments on the actual derivation}
\begin{proof}
Use the complete clause and parent list in
Annex~\ref{annex:p7.cr.small-php}, not an unspecified
``pigeonhole proof.'' For each listed assignment, all variables
other than the two named true literals are false. The following
table gives its route from node 19 to an input.
\begin{center}\small
\begin{tabular}{@{}lll@{}}
\toprule
Unmatched pigeon & True variables & Route node indices \\
\midrule
1 & \(x_2,y_3\) & \(19,17,16,14,1\) \\
1 & \(y_2,x_3\) & \(19,18,13,12,10,1\) \\
2 & \(x_1,y_3\) & \(19,17,15,2\) \\
2 & \(y_1,x_3\) & \(19,18,13,11,2\) \\
3 & \(x_1,y_2\) & \(19,18,3\) \\
3 & \(y_1,x_2\) & \(19,18,3\) \\
\bottomrule
\end{tabular}
\end{center}
At node 19, the y3 pivot selects node 17 when y3 is true and
node 18 when it is false. At node 18, x3 selects node 13 when
true and row 3 when false. At node 17, x2 selects node 16 when
true and node 15 when false. At node 16, false x1 selects
node 14; at node 15, false y2 selects row 2. At node 13, true
y2 selects node 12 and false y2 selects node 11. At node 12,
false y1 selects node 10; at node 10, false x1 selects row 1.
At node 11, false x2 selects row 2. These are precisely the
pivot choices for each row of the table. Every indicated
clause is false under its row's full assignment.

The six assignments are all the critical assignments because
their count is \(3\cdot2!=6\). The endpoint is always the
unmatched row, with two assignments per input row. Inputs 1
and 2 are reused in the derivation and duplicated on unfolding,
giving the already proved 19-versus-21 occurrence counts.
The assignments remain six distinct full truth assignments;
copying an input occurrence creates no additional assignment.
This verifies the concrete counting/representation distinction.
\end{proof}
\end{annexitem}
